% offset-mechanism-panels.tex
%
% v1: the four columns of SHOWCASE_offset_mechanism_<record> stated as red-pen edits
%     to the baseline matrices, and nothing else. No dynamics, no closed forms, no
%     attribution argument. Companions:
%       offset-mechanism-equations.tex   the same two fixes written as equation edits
%                                        per row, plus the disjoint-rows argument
%       offset-mechanism-derivation.tex  the formal Delta M, K = 0 and the two closed forms
%     This file overlaps offset-mechanism-equations.tex on purpose: that one is organised
%     by fix, this one is organised by figure panel. Input at most one of them.
%
% THIS IS A FRAGMENT, not a standalone document. It uses \mhr from
% jan-augmentation-writeup.tex and is meant to be pulled in with
%     % offset-mechanism-panels.tex
%
% v1: the four columns of SHOWCASE_offset_mechanism_<record> stated as red-pen edits
%     to the baseline matrices, and nothing else. No dynamics, no closed forms, no
%     attribution argument. Companions:
%       offset-mechanism-equations.tex   the same two fixes written as equation edits
%                                        per row, plus the disjoint-rows argument
%       offset-mechanism-derivation.tex  the formal Delta M, K = 0 and the two closed forms
%     This file overlaps offset-mechanism-equations.tex on purpose: that one is organised
%     by fix, this one is organised by figure panel. Input at most one of them.
%
% THIS IS A FRAGMENT, not a standalone document. It uses \mhr from
% jan-augmentation-writeup.tex and is meant to be pulled in with
%     % offset-mechanism-panels.tex
%
% v1: the four columns of SHOWCASE_offset_mechanism_<record> stated as red-pen edits
%     to the baseline matrices, and nothing else. No dynamics, no closed forms, no
%     attribution argument. Companions:
%       offset-mechanism-equations.tex   the same two fixes written as equation edits
%                                        per row, plus the disjoint-rows argument
%       offset-mechanism-derivation.tex  the formal Delta M, K = 0 and the two closed forms
%     This file overlaps offset-mechanism-equations.tex on purpose: that one is organised
%     by fix, this one is organised by figure panel. Input at most one of them.
%
% THIS IS A FRAGMENT, not a standalone document. It uses \mhr from
% jan-augmentation-writeup.tex and is meant to be pulled in with
%     % offset-mechanism-panels.tex
%
% v1: the four columns of SHOWCASE_offset_mechanism_<record> stated as red-pen edits
%     to the baseline matrices, and nothing else. No dynamics, no closed forms, no
%     attribution argument. Companions:
%       offset-mechanism-equations.tex   the same two fixes written as equation edits
%                                        per row, plus the disjoint-rows argument
%       offset-mechanism-derivation.tex  the formal Delta M, K = 0 and the two closed forms
%     This file overlaps offset-mechanism-equations.tex on purpose: that one is organised
%     by fix, this one is organised by figure panel. Input at most one of them.
%
% THIS IS A FRAGMENT, not a standalone document. It uses \mhr from
% jan-augmentation-writeup.tex and is meant to be pulled in with
%     \input{offset-mechanism-panels}
% Preview it with offset-mechanism-preview.tex. Requires amsmath.
%
% Supports SHOWCASE_offset_mechanism_<record> in scripts/gantry/msd-offset/figures/.

\subsection{What each panel of the showcase figure changes}
\label{sec:offset-panels}

The four columns of the showcase figure are four versions of the same model, differing only
in the edits below. All four are driven by the same recorded input, and every curve is that
version's open-loop simulation error against the recorded output.

\paragraph{Panel 1, \textit{baseline}.}
Three coordinates $q = [X,\ \Theta,\ Y]^\top$ obeying $M\ddot q + C\dot q + K q = F$. Writing
$m_t = m_1{+}m_2{+}m_b{+}m_h$, $\;b = \tfrac{(m_1-m_2)L_b}{2}$ and
$J = J_b{+}J_h{+}\tfrac{(m_1+m_2)L_b^2}{4} + m_h d^2$,
%
\begin{equation}
  M_{\mathrm{base}} =
  \begin{bmatrix}
    m_t        & b - m_h Y     & 0      \\
    b - m_h Y  & J + m_h Y^2   & -m_h d \\
    0          & -m_h d        & m_h
  \end{bmatrix},
  \qquad
  K_{\mathrm{base}} = \operatorname{diag}\bigl(0,\ k_{b1}{+}k_{b2},\ 0\bigr),
  \label{eq:panel-baseline}
\end{equation}
%
with $C_{\mathrm{base}}$ the usual matrix in $c_{g1},c_{g2},c_{b1},c_{b2},c_y$.

\paragraph{Panel 2, \textit{+ centre-of-mass term}.}
Split the payload as $m_h = \mhr + m_a$ and place $m_a$ at $Y + L_0$ instead of at $Y$, that is,
%
\begin{equation}
  m_h Y \;\longrightarrow\; \mhr Y + m_a (Y{+}L_0),
  \qquad
  m_h Y^2 \;\longrightarrow\; \mhr Y^2 + m_a (Y{+}L_0)^2 .
  \label{eq:panel-substitution}
\end{equation}
%
Exactly two entries of $M$ move as a result, and nothing else in the model does:
%
\begin{equation}
  M_{12} \;\longrightarrow\; \underbrace{b - m_h Y}_{\text{baseline}} - \; m_a L_0 ,
  \qquad
  M_{22} \;\longrightarrow\; \underbrace{J + m_h Y^2}_{\text{baseline}}
          + \; m_a\bigl[2 Y L_0 + L_0^2\bigr] .
  \label{eq:panel-fix1}
\end{equation}
%
$C$ and $K$ are untouched, the model keeps three coordinates, and no state is added. The size of
the edit is a first moment of mass, $m_a L_0 = 0.101\,\mathrm{kg\,m}$, which over
$m_h = 10.1\,\mathrm{kg}$ is a $1\,\mathrm{cm}$ error in the assumed position of the payload's
centre of mass.

\paragraph{Panel 3, \textit{+ absorber}.}
Leave $M$, $C$ and $K$ as \eqref{eq:panel-baseline} has them and add a fourth coordinate
$\delta_a$, the absorber's displacement relative to the payload. Each matrix gains one row and
one column:
%
\begin{equation}
  M =
  \left[\begin{array}{ccc|c}
    \multicolumn{3}{c|}{\multirow{3}{*}{$M_{\mathrm{base}}$}} & 0      \\
    \multicolumn{3}{c|}{}                                     & -m_a d \\
    \multicolumn{3}{c|}{}                                     & m_a    \\ \hline
    0 & -m_a d & m_a                                          & m_a
  \end{array}\right],
  \quad
  C =
  \left[\begin{array}{ccc|c}
    \multicolumn{3}{c|}{\multirow{3}{*}{$C_{\mathrm{base}}$}} & 0 \\
    \multicolumn{3}{c|}{}                                     & 0 \\
    \multicolumn{3}{c|}{}                                     & 0 \\ \hline
    0 & 0 & 0                                                 & c_a
  \end{array}\right],
  \quad
  K =
  \left[\begin{array}{ccc|c}
    \multicolumn{3}{c|}{\multirow{3}{*}{$K_{\mathrm{base}}$}} & 0 \\
    \multicolumn{3}{c|}{}                                     & 0 \\
    \multicolumn{3}{c|}{}                                     & 0 \\ \hline
    0 & 0 & 0                                                 & k_a
  \end{array}\right].
  \label{eq:panel-fix2}
\end{equation}
%
So $M_{Y\delta_a} = M_{\delta_a\delta_a} = m_a$ and $M_{\Theta\delta_a} = -m_a d$, while
$M_{X\delta_a} = 0$: the absorber has no direct path to $X$. The input $F$ still enters only the
first three rows, so no external force reaches the absorber.

\paragraph{Panel 4, \textit{+ both}.}
Both edits at once, \eqref{eq:panel-fix1} and \eqref{eq:panel-fix2}. This is the true $8$-state
plant, so this panel also shows the truth model's own residual against the recorded data.

\paragraph{The two edits are of different kinds.}
%
\begin{center}
\begin{tabular}{@{}lll@{}}
  \textit{Panel} & \textit{What changed} & \textit{Size of the model} \\[0.6mm]
  + centre-of-mass term & two numbers in $M$ & $3\times 3$, unchanged \\
  + absorber            & one new coordinate, one new row and column & $4\times 4$ \\
\end{tabular}
\end{center}
%
Panel 2 corrects a \emph{parameter}; panel 3 adds a \emph{state}. Neither can do the other's job,
which is why each panel removes the offset on one axis and leaves the other one unchanged.

\paragraph{One caveat, for a reader who opens the code.}
The coefficient $M_{12}$ of the full plant carries $-m_a(L_0 + \delta_a)$, not $-m_a L_0$. Panel 3
sets $L_0 = 0$ and therefore retains $-m_a \delta_a$, so its upper-left $3\times 3$ block equals
$M_{\mathrm{base}}$ only up to terms in $\delta_a$. Since $\delta_a$ has an rms of
$22\,\mu\mathrm{m}$ against $L_0 = 0.10\,\mathrm{m}$, that residual is worth $0.16\,\%$ and panel
3 reproduces the baseline $X$ offset to the quoted $100.00\,\%$ regardless.

% Preview it with offset-mechanism-preview.tex. Requires amsmath.
%
% Supports SHOWCASE_offset_mechanism_<record> in scripts/gantry/msd-offset/figures/.

\subsection{What each panel of the showcase figure changes}
\label{sec:offset-panels}

The four columns of the showcase figure are four versions of the same model, differing only
in the edits below. All four are driven by the same recorded input, and every curve is that
version's open-loop simulation error against the recorded output.

\paragraph{Panel 1, \textit{baseline}.}
Three coordinates $q = [X,\ \Theta,\ Y]^\top$ obeying $M\ddot q + C\dot q + K q = F$. Writing
$m_t = m_1{+}m_2{+}m_b{+}m_h$, $\;b = \tfrac{(m_1-m_2)L_b}{2}$ and
$J = J_b{+}J_h{+}\tfrac{(m_1+m_2)L_b^2}{4} + m_h d^2$,
%
\begin{equation}
  M_{\mathrm{base}} =
  \begin{bmatrix}
    m_t        & b - m_h Y     & 0      \\
    b - m_h Y  & J + m_h Y^2   & -m_h d \\
    0          & -m_h d        & m_h
  \end{bmatrix},
  \qquad
  K_{\mathrm{base}} = \operatorname{diag}\bigl(0,\ k_{b1}{+}k_{b2},\ 0\bigr),
  \label{eq:panel-baseline}
\end{equation}
%
with $C_{\mathrm{base}}$ the usual matrix in $c_{g1},c_{g2},c_{b1},c_{b2},c_y$.

\paragraph{Panel 2, \textit{+ centre-of-mass term}.}
Split the payload as $m_h = \mhr + m_a$ and place $m_a$ at $Y + L_0$ instead of at $Y$, that is,
%
\begin{equation}
  m_h Y \;\longrightarrow\; \mhr Y + m_a (Y{+}L_0),
  \qquad
  m_h Y^2 \;\longrightarrow\; \mhr Y^2 + m_a (Y{+}L_0)^2 .
  \label{eq:panel-substitution}
\end{equation}
%
Exactly two entries of $M$ move as a result, and nothing else in the model does:
%
\begin{equation}
  M_{12} \;\longrightarrow\; \underbrace{b - m_h Y}_{\text{baseline}} - \; m_a L_0 ,
  \qquad
  M_{22} \;\longrightarrow\; \underbrace{J + m_h Y^2}_{\text{baseline}}
          + \; m_a\bigl[2 Y L_0 + L_0^2\bigr] .
  \label{eq:panel-fix1}
\end{equation}
%
$C$ and $K$ are untouched, the model keeps three coordinates, and no state is added. The size of
the edit is a first moment of mass, $m_a L_0 = 0.101\,\mathrm{kg\,m}$, which over
$m_h = 10.1\,\mathrm{kg}$ is a $1\,\mathrm{cm}$ error in the assumed position of the payload's
centre of mass.

\paragraph{Panel 3, \textit{+ absorber}.}
Leave $M$, $C$ and $K$ as \eqref{eq:panel-baseline} has them and add a fourth coordinate
$\delta_a$, the absorber's displacement relative to the payload. Each matrix gains one row and
one column:
%
\begin{equation}
  M =
  \left[\begin{array}{ccc|c}
    \multicolumn{3}{c|}{\multirow{3}{*}{$M_{\mathrm{base}}$}} & 0      \\
    \multicolumn{3}{c|}{}                                     & -m_a d \\
    \multicolumn{3}{c|}{}                                     & m_a    \\ \hline
    0 & -m_a d & m_a                                          & m_a
  \end{array}\right],
  \quad
  C =
  \left[\begin{array}{ccc|c}
    \multicolumn{3}{c|}{\multirow{3}{*}{$C_{\mathrm{base}}$}} & 0 \\
    \multicolumn{3}{c|}{}                                     & 0 \\
    \multicolumn{3}{c|}{}                                     & 0 \\ \hline
    0 & 0 & 0                                                 & c_a
  \end{array}\right],
  \quad
  K =
  \left[\begin{array}{ccc|c}
    \multicolumn{3}{c|}{\multirow{3}{*}{$K_{\mathrm{base}}$}} & 0 \\
    \multicolumn{3}{c|}{}                                     & 0 \\
    \multicolumn{3}{c|}{}                                     & 0 \\ \hline
    0 & 0 & 0                                                 & k_a
  \end{array}\right].
  \label{eq:panel-fix2}
\end{equation}
%
So $M_{Y\delta_a} = M_{\delta_a\delta_a} = m_a$ and $M_{\Theta\delta_a} = -m_a d$, while
$M_{X\delta_a} = 0$: the absorber has no direct path to $X$. The input $F$ still enters only the
first three rows, so no external force reaches the absorber.

\paragraph{Panel 4, \textit{+ both}.}
Both edits at once, \eqref{eq:panel-fix1} and \eqref{eq:panel-fix2}. This is the true $8$-state
plant, so this panel also shows the truth model's own residual against the recorded data.

\paragraph{The two edits are of different kinds.}
%
\begin{center}
\begin{tabular}{@{}lll@{}}
  \textit{Panel} & \textit{What changed} & \textit{Size of the model} \\[0.6mm]
  + centre-of-mass term & two numbers in $M$ & $3\times 3$, unchanged \\
  + absorber            & one new coordinate, one new row and column & $4\times 4$ \\
\end{tabular}
\end{center}
%
Panel 2 corrects a \emph{parameter}; panel 3 adds a \emph{state}. Neither can do the other's job,
which is why each panel removes the offset on one axis and leaves the other one unchanged.

\paragraph{One caveat, for a reader who opens the code.}
The coefficient $M_{12}$ of the full plant carries $-m_a(L_0 + \delta_a)$, not $-m_a L_0$. Panel 3
sets $L_0 = 0$ and therefore retains $-m_a \delta_a$, so its upper-left $3\times 3$ block equals
$M_{\mathrm{base}}$ only up to terms in $\delta_a$. Since $\delta_a$ has an rms of
$22\,\mu\mathrm{m}$ against $L_0 = 0.10\,\mathrm{m}$, that residual is worth $0.16\,\%$ and panel
3 reproduces the baseline $X$ offset to the quoted $100.00\,\%$ regardless.

% Preview it with offset-mechanism-preview.tex. Requires amsmath.
%
% Supports SHOWCASE_offset_mechanism_<record> in scripts/gantry/msd-offset/figures/.

\subsection{What each panel of the showcase figure changes}
\label{sec:offset-panels}

The four columns of the showcase figure are four versions of the same model, differing only
in the edits below. All four are driven by the same recorded input, and every curve is that
version's open-loop simulation error against the recorded output.

\paragraph{Panel 1, \textit{baseline}.}
Three coordinates $q = [X,\ \Theta,\ Y]^\top$ obeying $M\ddot q + C\dot q + K q = F$. Writing
$m_t = m_1{+}m_2{+}m_b{+}m_h$, $\;b = \tfrac{(m_1-m_2)L_b}{2}$ and
$J = J_b{+}J_h{+}\tfrac{(m_1+m_2)L_b^2}{4} + m_h d^2$,
%
\begin{equation}
  M_{\mathrm{base}} =
  \begin{bmatrix}
    m_t        & b - m_h Y     & 0      \\
    b - m_h Y  & J + m_h Y^2   & -m_h d \\
    0          & -m_h d        & m_h
  \end{bmatrix},
  \qquad
  K_{\mathrm{base}} = \operatorname{diag}\bigl(0,\ k_{b1}{+}k_{b2},\ 0\bigr),
  \label{eq:panel-baseline}
\end{equation}
%
with $C_{\mathrm{base}}$ the usual matrix in $c_{g1},c_{g2},c_{b1},c_{b2},c_y$.

\paragraph{Panel 2, \textit{+ centre-of-mass term}.}
Split the payload as $m_h = \mhr + m_a$ and place $m_a$ at $Y + L_0$ instead of at $Y$, that is,
%
\begin{equation}
  m_h Y \;\longrightarrow\; \mhr Y + m_a (Y{+}L_0),
  \qquad
  m_h Y^2 \;\longrightarrow\; \mhr Y^2 + m_a (Y{+}L_0)^2 .
  \label{eq:panel-substitution}
\end{equation}
%
Exactly two entries of $M$ move as a result, and nothing else in the model does:
%
\begin{equation}
  M_{12} \;\longrightarrow\; \underbrace{b - m_h Y}_{\text{baseline}} - \; m_a L_0 ,
  \qquad
  M_{22} \;\longrightarrow\; \underbrace{J + m_h Y^2}_{\text{baseline}}
          + \; m_a\bigl[2 Y L_0 + L_0^2\bigr] .
  \label{eq:panel-fix1}
\end{equation}
%
$C$ and $K$ are untouched, the model keeps three coordinates, and no state is added. The size of
the edit is a first moment of mass, $m_a L_0 = 0.101\,\mathrm{kg\,m}$, which over
$m_h = 10.1\,\mathrm{kg}$ is a $1\,\mathrm{cm}$ error in the assumed position of the payload's
centre of mass.

\paragraph{Panel 3, \textit{+ absorber}.}
Leave $M$, $C$ and $K$ as \eqref{eq:panel-baseline} has them and add a fourth coordinate
$\delta_a$, the absorber's displacement relative to the payload. Each matrix gains one row and
one column:
%
\begin{equation}
  M =
  \left[\begin{array}{ccc|c}
    \multicolumn{3}{c|}{\multirow{3}{*}{$M_{\mathrm{base}}$}} & 0      \\
    \multicolumn{3}{c|}{}                                     & -m_a d \\
    \multicolumn{3}{c|}{}                                     & m_a    \\ \hline
    0 & -m_a d & m_a                                          & m_a
  \end{array}\right],
  \quad
  C =
  \left[\begin{array}{ccc|c}
    \multicolumn{3}{c|}{\multirow{3}{*}{$C_{\mathrm{base}}$}} & 0 \\
    \multicolumn{3}{c|}{}                                     & 0 \\
    \multicolumn{3}{c|}{}                                     & 0 \\ \hline
    0 & 0 & 0                                                 & c_a
  \end{array}\right],
  \quad
  K =
  \left[\begin{array}{ccc|c}
    \multicolumn{3}{c|}{\multirow{3}{*}{$K_{\mathrm{base}}$}} & 0 \\
    \multicolumn{3}{c|}{}                                     & 0 \\
    \multicolumn{3}{c|}{}                                     & 0 \\ \hline
    0 & 0 & 0                                                 & k_a
  \end{array}\right].
  \label{eq:panel-fix2}
\end{equation}
%
So $M_{Y\delta_a} = M_{\delta_a\delta_a} = m_a$ and $M_{\Theta\delta_a} = -m_a d$, while
$M_{X\delta_a} = 0$: the absorber has no direct path to $X$. The input $F$ still enters only the
first three rows, so no external force reaches the absorber.

\paragraph{Panel 4, \textit{+ both}.}
Both edits at once, \eqref{eq:panel-fix1} and \eqref{eq:panel-fix2}. This is the true $8$-state
plant, so this panel also shows the truth model's own residual against the recorded data.

\paragraph{The two edits are of different kinds.}
%
\begin{center}
\begin{tabular}{@{}lll@{}}
  \textit{Panel} & \textit{What changed} & \textit{Size of the model} \\[0.6mm]
  + centre-of-mass term & two numbers in $M$ & $3\times 3$, unchanged \\
  + absorber            & one new coordinate, one new row and column & $4\times 4$ \\
\end{tabular}
\end{center}
%
Panel 2 corrects a \emph{parameter}; panel 3 adds a \emph{state}. Neither can do the other's job,
which is why each panel removes the offset on one axis and leaves the other one unchanged.

\paragraph{One caveat, for a reader who opens the code.}
The coefficient $M_{12}$ of the full plant carries $-m_a(L_0 + \delta_a)$, not $-m_a L_0$. Panel 3
sets $L_0 = 0$ and therefore retains $-m_a \delta_a$, so its upper-left $3\times 3$ block equals
$M_{\mathrm{base}}$ only up to terms in $\delta_a$. Since $\delta_a$ has an rms of
$22\,\mu\mathrm{m}$ against $L_0 = 0.10\,\mathrm{m}$, that residual is worth $0.16\,\%$ and panel
3 reproduces the baseline $X$ offset to the quoted $100.00\,\%$ regardless.

% Preview it with offset-mechanism-preview.tex. Requires amsmath.
%
% Supports SHOWCASE_offset_mechanism_<record> in scripts/gantry/msd-offset/figures/.

\subsection{What each panel of the showcase figure changes}
\label{sec:offset-panels}

The four columns of the showcase figure are four versions of the same model, differing only
in the edits below. All four are driven by the same recorded input, and every curve is that
version's open-loop simulation error against the recorded output.

\paragraph{Panel 1, \textit{baseline}.}
Three coordinates $q = [X,\ \Theta,\ Y]^\top$ obeying $M\ddot q + C\dot q + K q = F$. Writing
$m_t = m_1{+}m_2{+}m_b{+}m_h$, $\;b = \tfrac{(m_1-m_2)L_b}{2}$ and
$J = J_b{+}J_h{+}\tfrac{(m_1+m_2)L_b^2}{4} + m_h d^2$,
%
\begin{equation}
  M_{\mathrm{base}} =
  \begin{bmatrix}
    m_t        & b - m_h Y     & 0      \\
    b - m_h Y  & J + m_h Y^2   & -m_h d \\
    0          & -m_h d        & m_h
  \end{bmatrix},
  \qquad
  K_{\mathrm{base}} = \operatorname{diag}\bigl(0,\ k_{b1}{+}k_{b2},\ 0\bigr),
  \label{eq:panel-baseline}
\end{equation}
%
with $C_{\mathrm{base}}$ the usual matrix in $c_{g1},c_{g2},c_{b1},c_{b2},c_y$.

\paragraph{Panel 2, \textit{+ centre-of-mass term}.}
Split the payload as $m_h = \mhr + m_a$ and place $m_a$ at $Y + L_0$ instead of at $Y$, that is,
%
\begin{equation}
  m_h Y \;\longrightarrow\; \mhr Y + m_a (Y{+}L_0),
  \qquad
  m_h Y^2 \;\longrightarrow\; \mhr Y^2 + m_a (Y{+}L_0)^2 .
  \label{eq:panel-substitution}
\end{equation}
%
Exactly two entries of $M$ move as a result, and nothing else in the model does:
%
\begin{equation}
  M_{12} \;\longrightarrow\; \underbrace{b - m_h Y}_{\text{baseline}} - \; m_a L_0 ,
  \qquad
  M_{22} \;\longrightarrow\; \underbrace{J + m_h Y^2}_{\text{baseline}}
          + \; m_a\bigl[2 Y L_0 + L_0^2\bigr] .
  \label{eq:panel-fix1}
\end{equation}
%
$C$ and $K$ are untouched, the model keeps three coordinates, and no state is added. The size of
the edit is a first moment of mass, $m_a L_0 = 0.101\,\mathrm{kg\,m}$, which over
$m_h = 10.1\,\mathrm{kg}$ is a $1\,\mathrm{cm}$ error in the assumed position of the payload's
centre of mass.

\paragraph{Panel 3, \textit{+ absorber}.}
Leave $M$, $C$ and $K$ as \eqref{eq:panel-baseline} has them and add a fourth coordinate
$\delta_a$, the absorber's displacement relative to the payload. Each matrix gains one row and
one column:
%
\begin{equation}
  M =
  \left[\begin{array}{ccc|c}
    \multicolumn{3}{c|}{\multirow{3}{*}{$M_{\mathrm{base}}$}} & 0      \\
    \multicolumn{3}{c|}{}                                     & -m_a d \\
    \multicolumn{3}{c|}{}                                     & m_a    \\ \hline
    0 & -m_a d & m_a                                          & m_a
  \end{array}\right],
  \quad
  C =
  \left[\begin{array}{ccc|c}
    \multicolumn{3}{c|}{\multirow{3}{*}{$C_{\mathrm{base}}$}} & 0 \\
    \multicolumn{3}{c|}{}                                     & 0 \\
    \multicolumn{3}{c|}{}                                     & 0 \\ \hline
    0 & 0 & 0                                                 & c_a
  \end{array}\right],
  \quad
  K =
  \left[\begin{array}{ccc|c}
    \multicolumn{3}{c|}{\multirow{3}{*}{$K_{\mathrm{base}}$}} & 0 \\
    \multicolumn{3}{c|}{}                                     & 0 \\
    \multicolumn{3}{c|}{}                                     & 0 \\ \hline
    0 & 0 & 0                                                 & k_a
  \end{array}\right].
  \label{eq:panel-fix2}
\end{equation}
%
So $M_{Y\delta_a} = M_{\delta_a\delta_a} = m_a$ and $M_{\Theta\delta_a} = -m_a d$, while
$M_{X\delta_a} = 0$: the absorber has no direct path to $X$. The input $F$ still enters only the
first three rows, so no external force reaches the absorber.

\paragraph{Panel 4, \textit{+ both}.}
Both edits at once, \eqref{eq:panel-fix1} and \eqref{eq:panel-fix2}. This is the true $8$-state
plant, so this panel also shows the truth model's own residual against the recorded data.

\paragraph{The two edits are of different kinds.}
%
\begin{center}
\begin{tabular}{@{}lll@{}}
  \textit{Panel} & \textit{What changed} & \textit{Size of the model} \\[0.6mm]
  + centre-of-mass term & two numbers in $M$ & $3\times 3$, unchanged \\
  + absorber            & one new coordinate, one new row and column & $4\times 4$ \\
\end{tabular}
\end{center}
%
Panel 2 corrects a \emph{parameter}; panel 3 adds a \emph{state}. Neither can do the other's job,
which is why each panel removes the offset on one axis and leaves the other one unchanged.

\paragraph{One caveat, for a reader who opens the code.}
The coefficient $M_{12}$ of the full plant carries $-m_a(L_0 + \delta_a)$, not $-m_a L_0$. Panel 3
sets $L_0 = 0$ and therefore retains $-m_a \delta_a$, so its upper-left $3\times 3$ block equals
$M_{\mathrm{base}}$ only up to terms in $\delta_a$. Since $\delta_a$ has an rms of
$22\,\mu\mathrm{m}$ against $L_0 = 0.10\,\mathrm{m}$, that residual is worth $0.16\,\%$ and panel
3 reproduces the baseline $X$ offset to the quoted $100.00\,\%$ regardless.
